% sections/03_model.tex
% Redesigned model section: JME register, airy exposition. Every block is motivated,
% every key equation interpreted. New devices: within-period timing enumeration, the
% price-system table (tab:prices), and the marginal adopter's indifference condition
% (from Section 5.2 of the working version) presented as the economics of the margin.
% Uses the preamble macros \DG \dsw \PEB \PEG \pE \coh \E \I \num.
% Labels provided for the rest of the paper: sec:model sec:firms sec:open sec:bop
% sec:ss eq:innernest. External refs: sec:cal sec:taste_id sec:experiments
% sec:conclusion sec:deflator sec:monetary app:booking app:taste.

\section{The model}\label{sec:model}

The economy is a small open economy that imports its fossil energy. It is populated
by heterogeneous households, competitive firms producing a domestic good from
labour, a union that sets wages, a government, and a rest of the world that sets
world prices. With one exception, the environment is the small open economy HANK of
\citet{ARS2024}: households consume a CES bundle of energy and goods, import prices
are sticky, monopoly profits are capitalised into a mutual fund, and monetary policy
holds the real interest rate constant.

The exception is the household's energy technology. In \citet{ARS2024} every
household buys energy at the same price with a fixed technology. Here each household
operates a discrete durable, brown or green, a gas boiler or a heat pump, a
combustion or an electric vehicle, and can pay a one-time cost $\psi_g$ to switch
from brown to green, after which it buys energy at the green price and is insulated
from the fossil price. This single change ramifies in three places, and only three:
the household block acquires a discrete technology state and an adoption choice;
the price system acquires a second household energy price and a technology-specific
cost of living; and the balance of payments acquires two flows, the switching cost
and the adopters' energy saving. Everything else is inherited.

Section~\ref{sec:hh} presents the household block, where the new economics lives.
Sections~\ref{sec:firms}--\ref{sec:gov} lay out the production side, the price
system, the open-economy accounts, and the government, which are standard.
Section~\ref{sec:eqm} defines equilibrium, and Section~\ref{sec:ss} the steady
state and the two crisis experiments. Time is discrete and quarterly. A bar denotes
a steady-state value; $x_{-1}$ and $x_{+1}$ a lag and an expected lead. All retail
prices are expressed relative to the CPI $P$ and carry the subscript ${}_{/P}$.

\subsection{Households}\label{sec:hh}

There are three permanent discount-factor types $i\in\{0,1,2\}$ of equal mass, with
discount factors $\beta_i$ chosen in Section~\ref{sec:cal} to match aggregate MPCs.
A household is described by three state variables: its assets $a$, its idiosyncratic
labour productivity $e$, and the durable technology $d\in\{B, G\}$ it operates.
Productivity follows a discretised log-AR(1), $\log e'=\rho_e\log e+\epsilon$, with a
mean-one grid and transition matrix $\Pi$, as in \citet{ARS2024}. The technology is
the new object, and because the paper concerns a decision over it, we begin by fixing
precisely when, within a period, the household holds it, loses it, and chooses it.

\subsubsection{The technology and the timing of the choice}\label{sec:tech}

A household enters the period operating the technology it chose last period, which we
denote $d_{-1}$. Three events then occur, in this order, before it consumes.

\begin{enumerate}
\item \emph{Breakdown.} A green durable fails with probability $\delta_g$ and reverts
to brown; a brown durable is unaffected. We write $d_-$ for the technology the
household holds \emph{after} this draw: the technology it brings to the decision.
Hence $d_-=B$ whenever $d_{-1}=B$, while $d_-=G$ with probability $1-\delta_g$ when
$d_{-1}=G$.
\item \emph{Productivity.} The household draws its labour productivity $e$.
\item \emph{The decision.} The household chooses, in a single problem, the
technology $d$ it will operate together with its consumption and saving. A household
holding brown ($d_-=B$) may keep brown or pay the switching cost $\psi_g$ to adopt
green; a household holding green ($d_-=G$) operates it, since running an owned green
durable costs nothing and reverting to brown is never chosen.
\end{enumerate}

No information arrives between the productivity draw and the decision, so the
household settles technology, consumption, and saving \emph{jointly}
(Section~\ref{sec:problem}). The technology it operates is what it carries
forward, $d$ becomes next period's $d_{-1}$, so the within-period sequence is
%
\[
\underbrace{d_{-1}}_{\text{held on entry}}
\;\xrightarrow{\ \text{breakdown}\ }\;
\underbrace{d_-}_{\text{held at the decision}}
\;\xrightarrow{\ \text{choice}\ }\;
\underbrace{d}_{\text{operated, carried forward}}.
\]
%
Two binary objects thus describe the technology: the holding at the decision, $d_-$,
and the operated technology, $d$. The switching cost is paid exactly when a brown
holder adopts, that is when $d_-=B$ and $d=G$; we write $\I[d_-=B, d=G]$ for a switcher
and $\I[d=G]$ for a household operating green. The feasible choices are
%
\[
\mathcal F(d_-)=
\begin{cases}
\{B,\,G\} & d_-=B \quad(\text{keep brown, or pay $\psi_g$ to adopt}),\\[2pt]
\{G\} & d_-=G \quad(\text{operate the green durable}).
\end{cases}
\]

Two features follow directly, and both recur below. First, the technology is
persistent but asymmetric: brown is absorbing, and green is lost only through
breakdown. Second, the adoption decision is, in effect, an option held by the brown
population alone: a green household has nothing to choose.

A constant green share therefore requires a standing flow of switchers to replace the
durables that fail. In steady state this flow is
%
\begin{equation}
\bar D^{\mathrm{sw}}=\delta_g\,\bar\DG,
\label{eq:switch-flow}
\end{equation}
%
with $\DG$ and $\dsw$ the population shares of green households and of switchers. At
the calibrated $\bar\DG=5\%$ and a five-year durable life, this flow is about a
quarter of one percent of households per quarter: adoption is a rare, lumpy event,
which is what makes its elasticity, not its level, the object to identify
(Section~\ref{sec:taste_id}, Appendix~\ref{app:taste}).

\subsubsection{Preferences and the cost of living}\label{sec:prefs}

Labour supply is determined at the union level (Section~\ref{sec:firms}), so flow
utility is CRRA over a consumption bundle $c$ alone,
%
\begin{equation}
u(c)=\frac{c^{\,1-\sigma}}{1-\sigma},
\qquad
u(c)=\log c \ \text{ for } \sigma=1,
\label{eq:felicity}
\end{equation}
%
with $\sigma$ the inverse elasticity of intertemporal substitution, as in \citet{ARS2024}. The bundle is the
\citet{ARS2024} CES nest of energy against a home/foreign goods composite, with
energy weight $\alpha_E$ and a low substitution elasticity $\eta_E$.

The departure from AMRS is that energy comes at a technology-specific price: a
household operating brown buys energy at $\PEB$, one operating green at $\PEG<\PEB$.
The operated technology therefore fixes both the energy price and the exact
cost-of-living index,
%
\begin{equation}
p^E_d=\PEB+\I[d=G]\,\bigl(\PEG-\PEB\bigr),
\label{eq:pEd}
\end{equation}
\begin{equation}
p_d=\Bigl[1+\alpha_E\bigl((p^E_d)^{1-\eta_E}-(\PEB)^{1-\eta_E}\bigr)\Bigr]^{\frac{1}{1-\eta_E}},
\qquad
p_d=\bigl(p^E_d/\PEB\bigr)^{\alpha_E}\ \text{ for }\eta_E=1.
\label{eq:pd}
\end{equation}
%
Because the CPI anchors its energy leg on the brown price
(equation~\eqref{eq:innernest} below), a brown household's cost of living is
identically one, $p_B\equiv1$, and $p_G<1$ records exactly the cheaper energy an
adopter enjoys: the object the switching cost buys.

\subsubsection{The household's problem}\label{sec:problem}

A household of type $i$ arrives at the consumption stage with assets $a$ and
productivity $e$, operating the technology $d$ it has just chosen, having held $d_-$
at the decision. Before writing the program it is worth stating, in words, what the
household has and what it does with it. On the resource side there are three items.
Assets earn the ex-post real return $r$, the return on the mutual fund of
Section~\ref{sec:firms}; it coincides with the lagged ex-ante rate
$r^{\mathrm{ante}}_{-1}$ at every date except the first, when the date-0 revaluation
of the fund's equity drives the two apart (equation~\eqref{eq:fund}), and the
household bears that surprise. Labour income is $e\,Z$, where $Z\equiv(1-\tau^L)wn$
is aggregate after-tax labour income distributed in proportion to productivity, as in
\citet{ARS2024}; the household does not choose hours, which the union sets
(Section~\ref{sec:firms}). The government may pay a transfer $T$, specified below. On
the spending side, the household buys the consumption bundle at its
technology-specific cost of living $p_d$, saves $a'\ge0$ subject to the borrowing
constraint, and, if it is a brown holder adopting green this period, pays the
switching cost $\psi_g$, once. The budget constraint collects these
items,\footnote{Formally the household solves in units of the domestic good,
the model's numeraire (Section~\ref{sec:eqm}): asset returns and wages are converted
at $p^{\num}=p_{H/P}$, and CPI-denominated objects, the transfers and $\psi_g$, are
divided by $p^{\num}$ on entry into \eqref{eq:coh}. We suppress the conversion factors
in the text.}
%
\begin{equation}
p_d\,c+a'+\psi_g\,\I[d_-=B,\,d=G]=(1+r)\,a+e\,Z+T.
\label{eq:coh}
\end{equation}

The switching cost is the economics of the margin in one term. It is paid exactly
once, in the period a brown holder adopts, and it buys the permanently cheaper
bundle $p_G<1$ for as long as the durable survives: an upfront, lumpy outlay set
against a gradual operating saving. Because the outlay is financed out of current
resources against the borrowing constraint, wealth governs who can undertake it,
which is why adoption concentrates among the unconstrained
(Section~\ref{sec:ss_adoption}).

The transfer carries ex-post fiscal policy,
%
\begin{equation}
T=T^{\mathrm{unt}}+ins_E\,(\pE-\bar p^E)\,\bar c^{\,B}_{E,i}(e, a),
\label{eq:transfer-hh}
\end{equation}
%
an untargeted lump sum $T^{\mathrm{unt}}$ plus a targeted (Slutsky) transfer that
scales, with insurance proportion $ins_E$, the crisis rise in the energy price by the
household's
\emph{pre-crisis} brown-energy demand $\bar c^{\,B}_{E,i}(e, a)$; both are zero in the
baseline. One design feature matters throughout: the transfer conditions on
steady-state energy use, never on the current technology. A household that switches
forfeits nothing, so the transfer leaves the adoption margin undistorted, in sharp
contrast to the price cap of Section~\ref{sec:prices}, which acts on $\PEB$ itself.

The household chooses its technology, consumption, and saving \emph{together}, in a
single problem:
%
\begin{equation}
\max_{d\in\mathcal F(d_-)}\;\;\max_{c,\;a'\ge0}\;
\Bigl\{\,u(c)+\beta_i\,\E\bigl[\bar V(d_-',e',a')\bigr]\Bigr\}
\quad\text{subject to \eqref{eq:coh}},
\label{eq:joint}
\end{equation}
%
with $\bar V$ the continuation value (defined in \eqref{eq:logit}) and next period's
holding $d_-'$ given by breakdown applied to the operated technology $d$. The nested
maxima are not a temporal sequence, since nothing is realised between them, but the
structure of a discrete--continuous choice: for each feasible technology the
household solves the continuous consumption-saving problem, obtaining the
technology-conditional value $V(d\mid d_-,e, a)$, in which the origin $d_-$ enters
only through the switching cost; it then selects the technology
(Section~\ref{sec:adoption}). Only three
continuous problems are ever relevant: a brown household ($p_B\equiv1$, no switching
cost), a green incumbent ($p_G$, no switching cost), and a green adopter ($p_G$,
resources reduced by $\psi_g$). Each satisfies the usual Euler and envelope
conditions,
%
\begin{equation}
u'(c)=\beta_i\,p_d\,\E\bigl[\bar V_a(d_-',e',a')\bigr],
\qquad
V_a=(1+r)\,u'(c)/p_d,
\label{eq:euler}
\end{equation}
%
with $u'(c)=c^{-\sigma}$; a green household's cheaper bundle raises the
consumption value of every unit of wealth. CES demands split by the operated
technology,
%
\begin{equation}
c_E(d)=\alpha_E\bigl(p^E_d/p_d\bigr)^{-\eta_E}c,
\qquad
c_{HF}(d)=(1-\alpha_E)\bigl(p_{HF/P}/p_d\bigr)^{-\eta_E}c,
\label{eq:hh-demands}
\end{equation}
%
and brown versus green energy use follow the technology,
$c^B_E=c_E\,(1-\I[d=G])$ and $c^G_E=c_E\,\I[d=G]$.

\subsubsection{The technology choice}\label{sec:adoption}

The outer maximum in \eqref{eq:joint} selects the technology, completing the joint
problem. A household holding $d_-$ compares the technology-conditional values
$V(d\mid d_-,\cdot)$; with an i.i.d.\ extreme-value taste shock of scale
$\sigma_\varepsilon$ on each option, the choice is logit and the household's
continuation value $\bar V$ is the log-sum-exp,
%
\begin{equation}
\begin{aligned}
\Pr(d\mid d_-,e, a)&=
\frac{\exp\{V(d\mid d_-,e, a)/\sigma_\varepsilon\}}
{\sum_{d'\in\mathcal F(d_-)}\exp\{V(d'\mid d_-,e, a)/\sigma_\varepsilon\}},\\[4pt]
\bar V(d_-,e, a)&=\sigma_\varepsilon\log\!\!\sum_{d'\in\mathcal F(d_-)}\!\!
\exp\bigl(V(d'\mid d_-,e, a)/\sigma_\varepsilon\bigr).
\end{aligned}
\label{eq:logit}
\end{equation}
%
For a green holder the feasible set is a singleton, so $\bar V(G,\cdot)=V(G\mid G,\cdot)$
and there is nothing to decide. The only genuine choice is the brown holder's: a
two-way comparison of staying brown against paying $\psi_g$ to adopt green. That cost
already sits inside the adopter's budget constraint \eqref{eq:coh}, so the option values
price it automatically.

The economics of the margin is transparent in steady state. The marginal switcher
equates the premium to the discounted operating-cost saving over the durable's
expected life,
%
\begin{equation}
\psi_g\;\approx\;\frac{\bigl(\bar P^E_B-\bar P^E_G\bigr)\,\bar C^G_E}{r+\delta_g},
\label{eq:margin}
\end{equation}
%
which ties the adoption threshold to the delivered cost ratio
$\varrho=\bar P^E_G/\bar P^E_B$. Every experiment in the paper acts on this one
comparison. An energy crisis raises $\PEB$ and lifts the return to switching; a
retail price cap holds $\PEB$ at its pre-crisis level and removes that return
entirely; a transfer supports income while leaving $\PEB$, and hence
\eqref{eq:margin},free; a permanent carbon tax raises the brown price in steady
state, greening the economy before any crisis arrives.

Away from steady state, the responsiveness of the margin is governed by
$\sigma_\varepsilon$. The taste shocks are, in the tradition of \citet{DCEGM2017},
the smoothing that makes the discrete choice differentiable; at a rare margin they are
also the structural semi-elasticity of switching, so $\sigma_\varepsilon$ cannot be
treated as a nuisance parameter and set small. Section~\ref{sec:taste_id} identifies
it from quasi-experimental adoption elasticities, and Appendix~\ref{app:taste}
develops the smoothing--elasticity duality.

\subsubsection{Aggregation}\label{sec:agg2}

For each type $i$, aggregates integrate the household policies over the stationary
distribution $\Lambda_i$, and the economy averages the three types equally,
%
\begin{equation}
X=\tfrac13\textstyle\sum_i X_i,
\qquad
X\in\{C,\;A,\;\DG,\;\dsw,\;C^B_E,\;C^G_E,\;C_{HF},\dots\}.
\label{eq:agg}
\end{equation}
%
Aggregate energy demand is $c_E=C^B_E+C^G_E$, and home- and foreign-good demands apply
the common inner nest to the composite,
$c_H=(1-\alpha_F)\,p_{H/HF}^{-\eta}\,C_{HF}$ and
$c_F=\alpha_F\,p_{F/HF}^{-\eta}\,C_{HF}$. The two aggregates the policy experiments
track are the green share $\DG$, the population share operating green, and the
switching flow $\dsw$, the share adopting green each period.

\subsection{Production and wage setting}\label{sec:firms}

\paragraph{Production and profits.}
Competitive firms produce the domestic final good from labour with productivity
$A^N$, sold at a constant gross markup $\mu$, so output is $y=A^N n$ and after-tax
aggregate labour income is $Z=(1-\tau^L)wn$, with $n$ hours and $w$ the real wage. Monopoly profits are capitalised into a mutual fund that households
hold,
%
\begin{equation}
D=(\mu-1)\,wn,
\qquad
J=D+\frac{J_{+1}}{1+r^{\mathrm{ante}}},
\qquad
j=\frac{J_{+1}}{1+r^{\mathrm{ante}}},
\qquad
1+r=\frac{J}{j_{-1}},
\label{eq:fund}
\end{equation}
%
following the notation of \citet{ARS2024}: $r^{\mathrm{ante}}_t$ is the ex-ante real
rate at which the risk-neutral fund prices its claims (equation~\eqref{eq:mp}), $j$ is
the ex-dividend share price, and $r$ is the fund's ex-post return: the return
households earn in \eqref{eq:coh}. Because the fund equates ex-ante returns across
assets, the ex-post return equals the lagged ex-ante rate at every date but the
first, $r_t=r^{\mathrm{ante}}_{t-1}$ for $t\ge1$; only the impact return $r_0$ differs,
reflecting the date-0 revaluation of the fund's equity, which is the source of the
revaluation term in \eqref{eq:nfa}. Importers are foreign-owned and the energy
endowment is not held domestically, so domestic firms are the only domestically held
risky claim.

\paragraph{Wage setting.}
A union sets the nominal wage under Calvo frictions $\theta_w$, trading off
households' marginal rate of substitution against the after-tax wage, with a
real-wage-stabilisation term (weight $\zeta_{BG}$) as in
\citet{ARS2024}. With slope $\kappa_w=(1-\theta_w)(1-\beta\theta_w)/\theta_w$,
inverse Frisch elasticity $\varphi$, and labour-disutility constant $v_\varphi$,
wage inflation follows
%
\begin{equation}
\pi^w(1+\pi^w)
= \beta\,\pi^w_{+1}(1+\pi^w_{+1})
+ \frac{\kappa_w}{v_\varphi\,n^{\varphi}}
\biggl[v_\varphi\,n^{\varphi}
- \frac{(1-\tau^L)w\,C^{-\sigma}}{\mu}
- \zeta_{BG}\frac{(w-\bar w)\,w\,\bar C^{-\sigma}\bar n}{\mu\,n\,\bar w}\biggr],
\label{eq:wpc}
\end{equation}
%
and the real wage evolves as $w=(1+\pi^w)\,w_{-1}/(1+\pi)$.

\subsection{The price system}\label{sec:prices}

The paper's mechanism runs entirely through relative prices, so it is worth laying
the price system out in one place. Table~\ref{tab:prices} collects it: world
prices at the top, a wholesale layer where nominal rigidities live, a retail layer
where policy operates, and the CPI that aggregates them.

\begin{table}[H]
\centering\small
\begin{tabular}{lll}
\toprule
Symbol & Object & Determined by \\
\midrule
$P^{E*},\,P^{*}_F,\,C^{*},\,r^{*}$ & world energy and goods prices, demand, rate &
exogenous\textsuperscript{a} \\
$\pE$ & wholesale retail price of brown energy & NKPC \eqref{eq:pE-nkpc} \\
$\PEB$ & brown households' energy price & cap rule \eqref{eq:cap} \\
$\PEG$ & green households' energy price & $\varrho\,\bar p^E$, eq.~\eqref{eq:pEG2} \\
$p_{H/P}$ & domestic good & markup pricing \eqref{eq:pH} \\
$p_{F/P}$ & imported consumption good & NKPC \eqref{eq:pF} \\
$p_{HF/P},\,p_{H/HF}$ & goods composites & CPI identities
\eqref{eq:innernest}--\eqref{eq:outernest} \\
$p_d$ & cost of living of technology $d$ & eq.~\eqref{eq:pd} \\
$Q$ & real exchange rate & UIP \eqref{eq:uip} \\
\bottomrule
\end{tabular}
\caption{The price system. All retail prices are relative to the CPI.
\textsuperscript{a}Under the fixed-quantity closure the world energy price
$P^{E*}$ clears the energy market (Section~\ref{sec:eqm}).}
\label{tab:prices}
\end{table}

\paragraph{The domestic good.}
Marginal-cost pricing at the constant markup gives
%
\begin{equation}
p_{H/P}=\frac{\mu}{A^N}\,w.
\label{eq:pH}
\end{equation}

\paragraph{Imported goods and wholesale energy.}
Retailers of the imported good and of energy reset prices under Calvo frictions
$\theta_F$ and $\theta_E$, so world price movements pass through to domestic
retail prices only gradually. With the usual slopes $\kappa_F$ and $\kappa_E$,
%
\begin{equation}
\kappa_F\Bigl(\frac{Q P^{*}_F}{p_{F/P}}-1\Bigr)
+ \beta_F\,\pi^F_{+1}(1+\pi^F_{+1}) - \pi^F(1+\pi^F)=0,
\label{eq:pF}
\end{equation}
\begin{equation}
\kappa_E\Bigl(\frac{Q P^{E*}}{\pE}-1\Bigr)
+ \beta_E\,\pi^E_{+1}(1+\pi^E_{+1}) - \pi^E(1+\pi^E)=0.
\label{eq:pE-nkpc}
\end{equation}
%
The margins these retailers earn out of steady state accrue abroad (the importers
are foreign-owned, as in \citealp{ARS2024}); they matter below only because the
import bill is valued at retail rather than world prices.

\paragraph{The retail energy prices, where policy lives.}
The price a brown household actually pays applies the cap to the wholesale price,
%
\begin{equation}
\PEB=(1-\tau^E)\,\pE+\tau^E\,\bar p^E,
\label{eq:cap}
\end{equation}
%
so $\tau^E$ dials pass-through continuously: at $\tau^E=0$ households face the
wholesale price, at $\tau^E=1$ the brown retail price is frozen at its pre-crisis
level: the full shield. The green price is fixed relative to the pre-crisis brown
price,
%
\begin{equation}
\PEG=\varrho\,\bar p^E,
\qquad \varrho<1,
\label{eq:pEG2}
\end{equation}
%
so adopters are fully insulated from the fossil shock.

\paragraph{The consumer price index.}
The CPI anchors its energy leg on the brown price alone: the statistical agency
prices the incumbent technology's bundle. The inner (energy versus goods) and
outer (home versus foreign) CES identities defining the composite prices are
%
\begin{equation}
(1-\alpha_E)\,p_{HF/P}^{\,1-\eta_E}+\alpha_E\,(\PEB)^{1-\eta_E}=1,
\label{eq:innernest}
\end{equation}
\begin{equation}
(1-\alpha_F)\,p_{H/HF}^{\,1-\eta}+\alpha_F\,p_{F/HF}^{\,1-\eta}=1.
\label{eq:outernest}
\end{equation}
%
Once green energy is cheaper and adopted, this index overstates the true cost of
living; the wedge has a closed form, closes no loop in the model, and is
quantified in Appendix~\ref{sec:deflator}.

\subsection{The open economy}\label{sec:open}\label{sec:bop}

\paragraph{Foreign demand and the exchange rate.}
Foreign demand for the home good is $c^{*}_H=\alpha^{*}(p^{*}_H)^{-\gamma}C^{*}$,
with trade elasticity $\gamma$, and the real exchange rate is pinned by uncovered
interest parity,
%
\begin{equation}
\frac{Q}{Q_{+1}}\bigl(1+r^{\mathrm{ante}}\bigr)=1+r^{*}.
\label{eq:uip}
\end{equation}

\paragraph{World energy supply.}
World supply is a fixed endowment $E^{*}$ plus intertemporal stock
arbitrage with cost $\Gamma$,
%
\begin{equation}
E^{\mathrm{stock}}=\frac{P^{E*}_{+1}/(1+r^{*})-P^{E*}}{\Gamma},
\qquad
E^{s}=E^{*}
+\bigl(E^{\mathrm{stock}}_{-1}-E^{\mathrm{stock}}\bigr).
\label{eq:energy-supply}
\end{equation}
%
The home economy owns none of the endowment, it is a pure importer, so energy
rents never enter domestic wealth.

\paragraph{The balance of payments and the adoption flows.}
The adoption margin adds exactly two terms to the current account, and they pull
in opposite directions. The switching cost is a resource outflow, booked as an
import $M^{\mathrm{sw}}$: it is front-loaded, paid once per switcher, and scales with
the switching volume. The adopters' saved fossil energy is a resource inflow, booked
as an import saving $S^{E}$: it accrues gradually, over the life of each green
durable,
%
\begin{equation}
M^{\mathrm{sw}}=\psi_g\,\dsw,
\qquad
S^{E}=(\PEB-\PEG)\,C^G_E.
\label{eq:adoption-flows}
\end{equation}
%
Net exports collect trade in goods and energy together with these two flows,
%
\begin{equation}
nx = Q\,p^{*}_H c^{*}_H - p_{F/P}\,c_F - \pE\,c_E
- M^{\mathrm{sw}} + S^{E},
\label{eq:nx}
\end{equation}
%
and net foreign assets accumulate at the ex-ante rate, with one additional term,
%
\begin{equation}
nfa = nx + (r-r^{\mathrm{ante}}_{-1})\bigl(A^{\mathrm{CPI}}_{-1}-j_{-1}\bigr)
+ (1+r^{\mathrm{ante}}_{-1})\,nfa_{-1}.
\label{eq:nfa}
\end{equation}
%
The middle term is a revaluation, and it is worth unpacking. Households earn the
ex-post return $r$ on their wealth, while the country's accounts accrue at the
promised ex-ante rate $r^{\mathrm{ante}}_{-1}$; the two differ only at date~0, when
the shock revalues the equity of the mutual fund (equation~\eqref{eq:fund}).
Applying the surprise $r-r^{\mathrm{ante}}_{-1}$ to households' wealth net of the
domestically held fund, $A^{\mathrm{CPI}}_{-1}-j_{-1}$, records the one-off capital
gain or loss on the country's external claims. At date~0, the only date the term is
active, $A^{\mathrm{CPI}}_{-1}-j_{-1}$ equals the net foreign position itself, since
government debt is zero at the pre-crisis steady state; from date~1 onward
$r_t=r^{\mathrm{ante}}_{t-1}$ and the term vanishes.

Booking the switching cost and the green energy as imports is a convention, not a
result: it fixes the sign of the adoption channel's \emph{output} contribution but
none of the policy rankings, which run through relative prices alone. The
alternative domestic booking is collected in Appendix~\ref{app:booking}.

\subsection{The government}\label{sec:gov}

\paragraph{Monetary policy.}
The nominal rate follows a rule in current and expected inflation, with the
real rate defined by the Fisher equation,
%
\begin{equation}
1+i=(1+\bar r^{*})(1+\phi_\pi\pi)(1+\phi_{\pi e}\pi_{+1})+\varepsilon^m,
\qquad
r^{\mathrm{ante}}=\frac{1+i}{1+\pi_{+1}}-1,
\label{eq:mp}
\end{equation}
%
where $\varepsilon^m$ is a monetary shock, zero outside the monetary
experiment of Appendix~\ref{sec:monetary}. The baseline is the constant-real-rate
rule of \citet{ARS2024}, $(\phi_\pi,\phi_{\pi e})=(0,1)$. Through \eqref{eq:uip} it
pins the path of the real exchange rate, so a domestic transfer has no
terms-of-trade effect: a property inherited from AMRS, not assumed anew. A standard
Taylor rule, $(\phi_\pi,\phi_{\pi e})=(1.5,0)$, is reported as a robustness
contrast in Appendix~\ref{sec:monetary}.

\paragraph{Fiscal policy.}
Fiscal policy is the object of Sections~\ref{sec:shield}--\ref{sec:exante}, so we lay
out each instrument explicitly. Four are crisis instruments, deployed once the shock
hits and zero in the baseline; a fifth, the carbon tax, is permanent and defines the
``prepared'' economy of Section~\ref{sec:exante}. Table~\ref{tab:instruments} states,
for each, what it acts on, the base its cost scales with, and its budgetary cost.

\begin{table}[H]
\centering\small
\begin{tabular}{llll}
\toprule
Instrument & Acts on & Cost scales with & Fiscal cost \\
\midrule
Price cap ($\tau^E$) & brown retail price $\PEB$ & actual energy use & $\tau^E(\pE-\bar p^E)\,c_E$ \\
Slutsky transfer ($ins_E$) & household income & pre-crisis brown use & $ins_E\,(\pE-\bar p^E)\,\bar C^B_E$ \\
Flat transfer ($T^{\mathrm{unt}}$) & household income & --- (lump sum) & $T^{\mathrm{unt}}$ \\
Green subsidy ($s_g$) & switching cost $\psi_g$ & number of switchers & $s_g\,\psi_g\,\dsw$ \\
\midrule
Carbon tax ($\tau_b$) & brown price $\PEB$, permanent & brown energy use & rebated (App.~\ref{app:booking}) \\
\bottomrule
\end{tabular}
\caption{Fiscal instruments. The first four are crisis instruments, zero in the
baseline (Sections~\ref{sec:shield}--\ref{sec:flat}); the carbon tax is permanent and
defines the prepared economy (Section~\ref{sec:exante}). Its revenue is rebated lump
sum in the baseline, with green-subsidy recycling in Appendix~\ref{app:booking}.}
\label{tab:instruments}
\end{table}

The distinction the results turn on is between acting on the \emph{price} and acting
on \emph{income}. The price cap holds $\PEB$ down through \eqref{eq:cap}; because
$\PEB$ is precisely what the adoption margin responds to (equation~\eqref{eq:margin}),
capping the price switches the margin off, and the cap's cost is a marginal
subsidy: it scales with actual crisis energy use. The two transfers instead leave
$\PEB$ free and support income, so the margin survives. They differ in their base: the
\emph{Slutsky} transfer (intensity $ins_E$) is indexed to each household's pre-crisis
brown-energy use, so a switcher keeps the same cheque and the margin is
undistorted, but, as Section~\ref{sec:flat} shows, indexing to energy use is
regressive; the \emph{flat} transfer $T^{\mathrm{unt}}$ is the equal-cost,
unconditional alternative. The \emph{green subsidy} pays a
fraction $s_g$ of the switching cost, forcing adoption directly rather than through
the price. The \emph{carbon tax} raises $\PEB$ permanently, greening the economy
before any crisis; it is the ex-ante counterpart to the crisis instruments and the
subject of Section~\ref{sec:exante}.

All crisis outlays are deficit-financed. Collecting them term by term, with each the
cost line of Table~\ref{tab:instruments},
%
\begin{equation}
X=\underbrace{T^{\mathrm{unt}}}_{\text{flat transfer}}
+\underbrace{ins_E\,(\pE-\bar p^E)\,\bar C^{B}_E}_{\text{Slutsky transfer}}
+\underbrace{\tau^E\,(\pE-\bar p^E)\,c_E}_{\text{price cap}}
+\underbrace{s_g\,\psi_g\,\dsw}_{\text{green subsidy}},
\label{eq:outlays}
\end{equation}
%
debt is stabilised by the distortionary labour tax,
%
\begin{equation}
B=(1+r^{\mathrm{ante}}_{-1})B_{-1}+X-\tau^L wn,
\qquad
\tau^L=\psi_B\,(B_{-1}-\bar B),
\label{eq:gbc2}
\end{equation}
%
with no debt in steady state ($\bar B=0$): crisis spending is repaid by raising the
labour tax in proportion to the debt it creates.

\subsection{Equilibrium}\label{sec:eqm}

\paragraph{Numeraire.}
The unit of account is the domestic good: households solve in domestic-good units
and their assets convert to CPI units as $A^{\mathrm{CPI}}=A\,p_{H/P}$. This
choice is not innocuous bookkeeping. With two household energy prices, any
CPI-based accounting would make the economy's resource constraint depend on how
the index weights brown and green energy; in domestic-good units the resource
constraint is invariant to that choice, and the CPI matters only where it
should, measurement and the monetary rule.

\paragraph{Market clearing.}
Goods, assets, and energy clear:
%
\begin{equation}
c_H+c^{*}_H-y=0,
\qquad
A^{\mathrm{CPI}}-nfa-j-B=0,
\qquad
\begin{cases}
P^{E*}\ \text{exogenous} & \text{(elastic supply)},\\
c_E=E^{s} & \text{(fixed quantity)},
\end{cases}
\label{eq:clearing}
\end{equation}
%
where the asset-market condition is redundant by Walras' law: it holds once goods,
energy, and the government budget balance, so we do not impose it separately. The two
energy closures correspond to the two experiments of Section~\ref{sec:experiments}: a
price-taking economy hit by an exogenous world-price path, and a fixed world quantity
whose price clears the market.

Let $\Psi_t$ denote the joint distribution over the household state
$(d, e, a)$, durable technology, idiosyncratic productivity, and assets, at date
$t$. Taking as given the sequences of prices, taxes, and transfers, each household
solves the recursive program of Section~\ref{sec:problem}: it chooses nondurable
consumption and savings by \eqref{eq:joint}--\eqref{eq:euler} and, at the durable
stage, switches from brown to green with the logit probability \eqref{eq:logit}.
This delivers value functions $V_t$ and policy functions
$\{c_t, a'_t, P^{sw}_t\}(d, e, a)$. The distribution then evolves as
$\Psi_{t+1}=\Lambda_t(\Psi_t)$, where the law of motion $\Lambda_t$ is induced by
these policies together with the exogenous productivity chain, and aggregates are the
corresponding moments of $\Psi_t$ \eqref{eq:agg}.

\paragraph{Definition (competitive equilibrium).}
Given the world paths $\{P^{E*},r^{*},P^{F*}\}$, the initial distribution $\Psi_0$
and net foreign assets, and the monetary and fiscal rules, a competitive equilibrium
is a sequence of aggregate prices
$\{w,\pi,\pi^{w},\pi^{F},\pi^{E},Q,\pE, r, r^{\mathrm{ante}},p_{H/P}\}_t$, value and
policy functions $\{V_t, c_t, a'_t, P^{sw}_t\}$, distributions $\{\Psi_t\}$, aggregate
quantities $\{y, n, c_H, c_F, c_E, C, A\}_t$, technology shares $\{\DG,\dsw\}_t$, and fiscal
variables $\{B,\tau^{L}\}_t$ such that:
\begin{enumerate}[label=(\roman*),leftmargin=2.2em, itemsep=1pt]
\item \emph{(Households)} given prices, taxes, and transfers, the value and policy
functions solve the recursive problem \eqref{eq:joint}--\eqref{eq:euler}, with
brown-to-green switching governed by \eqref{eq:logit}--\eqref{eq:margin};
\item \emph{(Consistency)} $\Psi_t$ evolves from $\Psi_0$ under the law of motion
induced by the policy functions and the productivity process, and the aggregates in
\eqref{eq:agg} are the moments of $\Psi_t$;
\item \emph{(Firms and unions)} the price and wage Phillips curves
\eqref{eq:fund}--\eqref{eq:wpc} and the price system
\eqref{eq:pH}--\eqref{eq:pEG2} hold;
\item \emph{(Government)} the monetary rule \eqref{eq:mp}, the fiscal rule, and the
government budget \eqref{eq:gbc2} hold;
\item \emph{(Open economy)} uncovered interest parity \eqref{eq:uip}, the
balance-of-payments and net-foreign-asset accounting
\eqref{eq:nx}--\eqref{eq:nfa}, and the energy-supply closure
\eqref{eq:energy-supply} hold;
\item \emph{(Clearing)} the goods and energy markets clear \eqref{eq:clearing}; the
asset market clears by Walras' law.
\end{enumerate}

This is the object of interest: a sequence of prices and allocations, together with a
distribution consistent with individual optimisation. The sequence-space Jacobian
method of Section~\ref{sec:ss} computes the linearised transition of this
equilibrium.

\subsection{Steady state and the two experiments}\label{sec:ss}

As a pure importer, the openness normalisations of \citet{ARS2024} reduce to
$\alpha_F=(\alpha-\alpha_E)/(1-\alpha_E)$ and a trade elasticity $\gamma=\eta$ set
to hit an import-demand target $\chi=0.3$, with scale normalisations $y=1$,
$A^N=\mu$, $\alpha^{*}=\alpha$, and $\bar E^{*}=\alpha_E$. The steady
state solves for the labour-disutility shifter $\phi$, the top discount factor
$\beta_{\max}$, output, and the switching cost $\psi_g$, against the wage Phillips
curve, the net-foreign-asset condition, goods clearing, and the green-share target
$\bar\DG=5\%$. The switching cost is thus calibrated to the level of adoption, and
Section~\ref{sec:cal} takes up the parameter the level cannot pin,
$\sigma_\varepsilon$. The steady state is identical across all baseline
experiments.

Two experiments drive the crisis, one per energy closure. The \emph{price} shock
doubles the world energy price on impact and lets it decay with a 16-quarter
half-life under perfectly elastic supply,
%
\begin{equation}
P^{E*}_t=\bar P^{E*}\bigl(1+\text{size}\cdot\rho^{\,t}\bigr),
\qquad
\rho=2^{-1/\text{half-life}},
\label{eq:price-shock}
\end{equation}
%
the experiment of \citet{ARS2024}. The \emph{supply} shock removes 10\% of the
world quantity for six quarters and lets the world price clear the market,
%
\begin{equation}
E^{*}_t=\bar E^{*}
\bigl(1-\text{drop}\cdot\I[t<\text{quarters}]\bigr),
\label{eq:supply-shock}
\end{equation}
%
the closure of \citet{Bayer2026}. Section~\ref{sec:experiments} puts the adoption
margin to work under both.
